\section{Introduction}\label{sec:introduction}
Encryption conceals message content, yet LLM engines and agents expose packet sizes, cache timing, visited domains and tool sequences \cite{weiss2024prompt,promptpeek2025,earlybird2025,agentprint2025,uitraces2026,jeong2026wra}. Defenses follow: padding, batching, cache isolation, decoy queries and action humanization. Each is evaluated at one receiver, execution boundary and workload, yet its headline number is cited as a property of the mechanism. A padding result measured on an upstream link says nothing about a frontend that re-emits decoded text; decoys whose read responses are discarded say nothing about cover that executes writes; and the overhead of a semantic cache does not price prefix isolation reported in the same paper.

We therefore systematize defense evidence at the level of the \emph{evaluated result record} rather than the paper. A record binds one outcome to its target, receiver projection, intervention boundary, attacker fit, workload and cost object (\S\ref{sec:scope}). Three instruments operate on these records: a map of 82 study families by signal origin, receiving view and threat polarity; a nine-item audit of what each mitigation family binds; and a six-axis matrix of reported holdouts. Together they answer one question: \emph{which changes invalidate a measured defense result?}

\paragraph{Contributions}
\begin{itemize}
\item \textbf{Record-level systematization (\S\S\ref{sec:methodology}--\ref{sec:discussion}).} Separating records exposes within-paper substitution: 7 of 15 same-family record pairs measure different primitives, targets or cost objects, and only 16 of 102 related record pairs support even a descriptive comparison. Reuse evidence is scarcest where deployment needs it: matched adaptation is bound for 2/11 applicable mitigation families and protected-point utility for 3/17, and no family reports an explicit hardware or temporal holdout.
\item \textbf{Four boundary demonstrations (\S\ref{sec:controlled-placement}).} Controlled experiments on a retail agent workload instantiate the emitter, executed-state, calibration and attacker-observation boundaries identified by the audit, each priced with its own cost object.
\item \textbf{Directional calibration transfer.} Moving calibration from MiniMax to GLM changes cancellation/return BA from 50.0/67.3\% to 70.8/50.0\%: protection moves in opposite directions, so no single transfer score summarizes it. Write-local shaping is inexpensive ($1.45$--$1.51\times$ bytes) but leaves return BA at 61.3--61.6\%, and a fresh live-planner study (37/50 versus 33/50 successes) bounds the utility claim of prefix cover.
\end{itemize}
Adaptive attackers, distribution shift and online cost evaluation are established website-fingerprinting (WF) practice \cite{juarez2014critical,mathews2023critical}; we apply them to records rather than claim them. Exps~1--2 deliberately reproduce known mechanisms so that their cost and state consequences can be measured; Exps~3--4 measure agent-specific operating points. Table~\ref{tab:rq-findings} links each research question to its finding.
\begin{table*}[t]\centering\footnotesize
\caption{Research questions, systematization findings and where each is established. Experiments test the comparison rules; the SoK contribution is the record-level synthesis.}\label{tab:rq-findings}
\begin{tabularx}{\textwidth}{@{}>{\raggedright\arraybackslash}p{.15\textwidth}Y>{\raggedright\arraybackslash}p{.29\textwidth}@{}}\toprule
Question & Finding & Evidence\\\midrule
RQ1: targets and observables & A common signal can serve confidentiality, attribution, evasion or authorized diagnostics; polarity changes the protected interest. Signal generation clusters at serving (32 families) and orchestration (19), whereas observation clusters at application, hardware and transport. & Complete 58-family core map (\S\ref{sec:corpus-map}, Table~\ref{tab:core-family-map}).\\
RQ2: origin versus access & Packet-size cover does not bind call order; plaintext re-emission creates a fresh size trace; shared call projections do not establish unchanged state. Within one paper, 7 of 15 same-family record pairs measure different primitives, targets or cost objects. & 39-record ledger (\S\ref{sec:defenses}); NetEcho/SplitDecoy and WRA/AHB lenses (\S\ref{sec:discussion}); Exps~1--2.\\
RQ3: portability & Explicit holdouts exist for task/input in 13 of 82 families, for model in three, and for neither hardware nor time in any; input separation does not certify a changed generator. & Six-axis matrix (\S\ref{sec:evaluation}, Appendix~\ref{app:holdouts}); calibration-transfer matrix in Exp~3.\\
RQ4: protection, state and cost & Inference is bound far more often than reuse: matched adaptation 2/11 applicable families, utility at the protected operating point 3/17, cost 11/19. & Nine-item, 20-family audit (\S\ref{sec:defenses}); Exps~3--4 leakage/state/cost contrasts.\\\bottomrule
\end{tabularx}\end{table*}


\soktakeaway{Scope of claims}{The real-agent target is a requested operation category in one retail environment with a known 53-customer cohort. Measurements establish fitted-suite inference, logical state conformance and scoped resource costs; a separate 100-episode live study measures prefix-cover utility under injected deadlines. Deployment-wide security and full joint utility require the validation specified in \S\ref{sec:validity}.}

Figure~\ref{fig:system-model} places trace and action paths on the same agent--tool chain, making the receiving representation and the execution boundary explicit before a defense is selected.
\begin{figure*}[t]
\centering\begin{tikzpicture}[font=\scriptsize,>=Latex,
box/.style={draw,rounded corners,align=center,text width=3.35cm,minimum height=10mm,fill=blue!5},
boundary/.style={box,fill=orange!12},
observer/.style={box,fill=green!8}]
\node[box] (llm) at (0,0) {LLM service\\token stream / cache activity};
\node[boundary] (api) at (4.25,0) {Upstream protected link\\padding or batching};
\node[boundary] (front) at (8.5,0) {Content-aware frontend\\decoded content $\rightarrow$ new stream};
\node[observer] (net) at (12.75,0) {Final network observer\\groups, lengths, timing};
\draw[->] (llm)--(api);
\draw[->] (api)--(front);
\draw[->] (front)--(net);
\node[box] (policy) at (0,-2.1) {Agent action policy\\task goal / current state};
\node[boundary] (cover) at (4.25,-2.1) {Cover boundary\\discard, shadow, or execute};
\node[box] (dispatch) at (8.5,-2.1) {Tool / GUI dispatch\\real and admitted cover actions};
\node[observer] (state) at (12.75,-2.1) {Environment / receiving site\\events and changed task state};
\draw[->] (policy)--(cover);
\draw[->] (cover)--(dispatch);
\draw[->] (dispatch)--(state);
\draw[->,dashed] (llm.south)--node[left,align=right]{action\\selection}(policy.north);
\node[font=\scriptsize,align=center] at (6.37,-.95)
{Trace path: upstream protection must be related to the final emitting interface (NetEcho).};
\node[font=\scriptsize,align=center] at (6.37,-3.0)
{Action path: executed cover can change state (AHB); timestamp-only edits can retain event structure (UI Traces).};
\end{tikzpicture}

\caption{Observation and execution boundaries. The four demonstrations test emitter placement, stateful execution/link composition, real-agent operation signatures, and timing with matched transport controls.}
\label{fig:system-model}
\end{figure*}
